\chapter{Audit et inventaire de l'existant}
\label{chap:audit}

\begin{todobox}[Chapitre à compléter après la visite de site]
Ce chapitre contient la méthodologie d'audit et la checklist à remplir. Les résultats sont à intégrer à l'issue de la visite prévue dans les deux semaines.
\end{todobox}

\section{Objectif de l'audit}
L'audit a pour objectif de produire un inventaire exhaustif et fiable de l'infrastructure physique, réseau et stockage existante, afin de déterminer :
\begin{itemize}
  \item Ce qui peut être réutilisé dans la plateforme cible.
  \item Ce qui doit être remplacé.
  \item Ce qui manque pour atteindre le dimensionnement cible.
\end{itemize}

\begin{warnbox}[Règle de sécurité]
Aucune action destructive ne doit être réalisée pendant l'audit. Ne pas démarrer, éteindre, débrancher ou modifier la configuration d'un équipement sans autorisation explicite et sans avoir identifié son rôle potentiel.
\end{warnbox}

\section{Checklist — visite de site}

\subsection{Inventaire physique — serveurs}
\begin{longtable}{@{}p{0.5cm}p{5cm}p{4.5cm}p{3.5cm}@{}}
\toprule
\textbf{\#} & \textbf{Action} & \textbf{Relevé} & \textbf{Commentaire} \\
\midrule
\endfirsthead
\toprule
\textbf{\#} & \textbf{Action (suite)} & \textbf{Relevé} & \textbf{Commentaire} \\
\midrule
\endhead
\midrule
\multicolumn{4}{r}{\footnotesize\itshape Suite page suivante \ldots} \\
\endfoot
\bottomrule
\endlastfoot
1 & Photographier chaque serveur (face avant, arrière, étiquette de service) & \todo & \\
2 & Relever marque, modèle, numéro de série & \todo & \\
3 & Identifier les CPUs (référence, nombre de sockets, cœurs) & \todo & \\
4 & Relever la quantité et le type de RAM installée & \todo & \\
5 & Identifier le contrôleur RAID/HBA (modèle, firmware) & \todo & \\
6 & Relever le nombre de baies et le type (2,5" / 3,5") & \todo & \\
7 & Identifier les disques installés (référence, taille, état SMART) & \todo & \\
8 & Relever les cartes réseau (nombre de ports, débit, fabricant) & \todo & \\
9 & Identifier la fonction de gestion (iDRAC, iLO) et l'adresse IP & \todo & \\
10 & Vérifier les alimentations (nombre, état, raccordement) & \todo & \\
11 & Vérifier si le serveur est allumé et son état (LED) & \todo & \\
12 & Vérifier si les données sont présentes ou si les disques sont vierges & \todo & \\
\end{longtable}

\subsection{Inventaire physique — stockage}
\begin{longtable}{@{}p{0.5cm}p{5cm}p{4.5cm}p{3.5cm}@{}}
\toprule
\textbf{\#} & \textbf{Action} & \textbf{Relevé} & \textbf{Commentaire} \\
\midrule
\endfirsthead
\toprule
\textbf{\#} & \textbf{Action (suite)} & \textbf{Relevé} & \textbf{Commentaire} \\
\midrule
\endhead
\midrule
\multicolumn{4}{r}{\footnotesize\itshape Suite page suivante \ldots} \\
\endfoot
\bottomrule
\endlastfoot
13 & Identifier les baies de stockage (marque, modèle) & \todo & \\
14 & Relever la connectique (SAS, FC, iSCSI, NFS) & \todo & \\
15 & Vérifier l'état (allumée, éteinte, LED) & \todo & \\
16 & Identifier les disques installés et leur état & \todo & \\
17 & Récupérer la configuration si possible (via console) & \todo & \\
18 & Vérifier si des données sont présentes & \todo & \\
\end{longtable}

\subsection{Inventaire physique — réseau}
\begin{longtable}{@{}p{0.5cm}p{5cm}p{4.5cm}p{3.5cm}@{}}
\toprule
\textbf{\#} & \textbf{Action} & \textbf{Relevé} & \textbf{Commentaire} \\
\midrule
\endfirsthead
\toprule
\textbf{\#} & \textbf{Action (suite)} & \textbf{Relevé} & \textbf{Commentaire} \\
\midrule
\endhead
\midrule
\multicolumn{4}{r}{\footnotesize\itshape Suite page suivante \ldots} \\
\endfoot
\bottomrule
\endlastfoot
19 & Identifier le firewall Juniper (modèle exact, version JunOS) & SRX300 (à confirmer) & \\
20 & Relever la configuration actuelle du SRX (si accessible) & \todo & \\
21 & Identifier les switches Cisco (modèle, nombre de ports, débit) & \todo & \\
22 & Vérifier si les switches sont administrables (SSH, console) & Oui (confirmé) & \\
23 & Relever la configuration actuelle (VLANs, trunks, LACP) & \todo & \\
24 & Identifier l'arrivée fibre et le raccordement opérateur & \todo & \\
25 & Vérifier le raccordement du bloc public /28 & SRX300 (confirmé) & \\
26 & Identifier les liens entre switches (topologie) & \todo & \\
27 & Vérifier la présence de câbles 10 GbE (SFP+) & \todo & \\
\end{longtable}

\subsection{Inventaire physique — environnement}
\begin{longtable}{@{}p{0.5cm}p{5cm}p{4.5cm}p{3.5cm}@{}}
\toprule
\textbf{\#} & \textbf{Action} & \textbf{Relevé} & \textbf{Commentaire} \\
\midrule
\endfirsthead
\toprule
\textbf{\#} & \textbf{Action (suite)} & \textbf{Relevé} & \textbf{Commentaire} \\
\midrule
\endhead
\midrule
\multicolumn{4}{r}{\footnotesize\itshape Suite page suivante \ldots} \\
\endfoot
\bottomrule
\endlastfoot
28 & Vérifier la présence et l'état d'un onduleur & Aucun connu & \\
29 & Vérifier la capacité de l'onduleur (kVA, autonomie) & \todo & \\
30 & Vérifier le nombre de circuits électriques disponibles & \todo & \\
31 & Vérifier la redondance des alimentations serveurs & \todo & \\
32 & Vérifier la climatisation (type, capacité, redondance) & Présente & \\
33 & Relever la température et l'hygrométrie de la salle & \todo & \\
34 & Identifier les accès physiques (serrures, badges) & \todo & \\
\end{longtable}

\subsection{Documentation à récupérer}
\begin{itemize}
  \item Schémas réseau de l'ancienne infrastructure.
  \item Inventaire matériel historique.
  \item Documentation OpenStack (version, configuration, projets).
  \item Configuration Ceph (si utilisée).
  \item Fichiers Ansible / scripts de déploiement.
  \item Documentation Proxmox (si utilisée).
  \item Plans de câblage.
  \item Contrats de support et licences en cours.
  \item Contacts des anciens prestataires.
\end{itemize}

\section{Résultats de l'audit}
\begin{todobox}[À compléter après la visite]
Cette section accueillera le tableau d'inventaire consolidé, les photographies, et la liste des éléments réutilisables, à remplacer et à acquérir.
\end{todobox}